\begin{corollary}[Chow tensor summands of hyper-K\"ahler varieties]
\label{c:hyperkahler-chow-unramified}
\status{Proved using the stated 6 October 2026 purity theorem}
Let $M$ be a pure rational Chow motive over a number field, obtained
from products of hyper-K\"ahler varieties by homogeneous algebraic
projectors, finite sums, duals and Tate twists.  At a finite place
$v$ of residue characteristic $p$, if $T_{\ell_0}(M)$ is unramified
for one $\ell_0\ne p$, then $T_\ell(M)$ is unramified for every
$\ell\ne p$.
\end{corollary}

\begin{proof}
\citet[Theorem~4.4, with Lemmas~4.2--4.3]{c_itozou2026}
prove weight--monodromy in every degree for every hyper-K\"ahler
variety over a $p$-adic field, including $\ell=p$.  Their theorem
is over the original local field: the finite extension in its proof
does not remain in its conclusion, because purity descends through
finite extensions.  Pure Weil--Deligne representations are stable
under tensor products, duals, Tate twists and direct summands.
K\"unneth therefore gives the ambient purity hypothesis in Proposition~\ref{c:chow-unramified-purity-reduction}.  At $\ell_0$, zero monodromy gives
a single Frobenius weight on the chosen pure summand.  That reduction,
using Saito's algebraic-correspondence trace theorem, proves the
assertion at the original place $v$.
\end{proof}

This consequence uses the purity theorem, not a descent of
unramifiedness through a ramified extension.  The stronger
group-valued compatibility statement in
\citet[Theorem~5.8 and Remark~5.9]{c_itozou2026} assumes connected
monodromy and a semistable Kuga--Satake variety over the chosen base;
Corollary~5.10 obtains its general local Frobenius comparison after
a finite extension.  The displayed corollary retains algebraic
projectors.  The same purity theorem gives the local C.4 bound
$\operatorname{weights}(\ker N)\le j$ for pure weight-$j$
tensor summands of hyper-K\"ahler cohomology.  It does not assert the missing trace comparison for
arbitrary absolute-Hodge projectors outside the abelian-type case.
